\subsection{Analiza matematyczna dla Algorytmu \ref{alg:wielowymiarowy} (k $\ge$ 1)}

W niniejszym podrozdziale przeprowadzimy analizę matematyczną uogólnionego algorytmu dla przypadku wielowymiarowego ($k \ge 1$). Głównym wyzwaniem, odróżniającym tę sytuację od wariantu jednowymiarowego, jest fakt, że poszukujemy całej podprzestrzeni rozpiętej przez macierz parametrów $A$, a nie tylko pojedynczego wektora . Wymaga to zastosowania rozkładu według wartości osobliwych (SVD) i zaawansowanych probabilistycznych nierówności macierzowych.

\subsubsection{Główne twierdzenie o błędzie aproksymacji}

Podobnie jak poprzednio, analizę rozpoczynamy od sformułowania ostatecznego ograniczenia na błąd aproksymacji.

\begin{tw}[Zbieżność Algorytmu \ref{alg:wielowymiarowy} \cite{fornasier2012learning}]\label{tw:zbieznoscalgorytmu2}
Niech $\log d \le m_\Phi \le [\log 6]^2 d$. Wówczas istnieje stała $c_1'$ taka, że wykorzystując $m_\mathcal{X} \cdot (m_\Phi + 1)$ ewaluacji funkcji $f$, Algorytm 2 definiuje funkcję $\hat{f}: B_{\mathbb{R}^d}(1+\bar{\epsilon}) \to \mathbb{R}$, która z prawdopodobieństwem co najmniej
$$1 - \left(e^{-c_1' m_\Phi} + e^{-\sqrt{m_\Phi d}} + k e^{\frac{-m_\mathcal{X} \alpha s^2}{2kC_2^2}}\right)$$
spełnia nierówność
$$\|f - \hat{f}\|_\infty \le 2C_2\sqrt{k}(1+\bar{\epsilon})\frac{\nu_2}{\sqrt{\alpha(1-s)} - \nu_2},$$
gdzie
$$\nu_2 = C\left(k^{1/q}\left[\frac{m_\Phi}{\log(d/m_\Phi)}\right]^{1/2 - 1/q} + \frac{\epsilon k^2}{\sqrt{m_\Phi}}\right),$$
a stała $C$ zależy tylko od $C_1$ i $C_2$ definiujących model .
\end{tw}

\subsubsection{Rekonstrukcja macierzy pomiarowej}

W pierwszym kroku musimy oszacować błąd odzyskania całej macierzy $X$ za pomocą optymalizacji $l_1$. Aplikując Twierdzenie \ref{tw:aproksymacjaK-czlonowa} kolumna po kolumnie, uzyskujemy następujący wniosek :

\begin{wn}\label{wn:bladodzyskaniaX}
Niech $\log d \le m_\Phi \le [\log 6]^2 d$. Wówczas z prawdopodobieństwem co najmniej $1 - (e^{-c_1' m_\Phi} + e^{-\sqrt{m_\Phi d}})$ macierz $\hat{X}$ wyliczona w Algorytmie 2 spełnia:
$$\|X - \hat{X}\|_F \le C\sqrt{m_\mathcal{X}}\left(k^{1/q}\left[\frac{m_\Phi}{\log(d/m_\Phi)}\right]^{1/2 - 1/q} + \frac{\epsilon k^2}{\sqrt{m_\Phi}}\right) = \sqrt{m_\mathcal{X}}\nu_2,$$
gdzie $C$ zależy wyłącznie od $C_1$ i $C_2$ .
\end{wn}

\begin{proof}
Dowód przebiega analogicznie do Wniosku \ref{wn:bladodzyskaniakolumn}. Rozkładamy normę Frobeniusa na sumę błędów poszczególnych kolumn: $\|X - \hat{X}\|_F^2 = \sum_{j=1}^{m_\mathcal{X}} \|x_j - \hat{x}_j\|_{l_2^d}^2$ . Najlepszą $K$-członową aproksymację dla kolumny szacujemy używając własności normy macierzy $A$: $\|x_j\|_{l_q^d} \le C_1 C_2 k^{1/q}$. Po wstawieniu tych wartości oraz oszacowaniu norm błędów Taylora $\|\epsilon_j\|_{l_2^{m_\Phi}}$ i $\|\epsilon_j\|_{l_\infty^{m_\Phi}}$ z uwzględnieniem czynnika $k^2$ , a następnie po zsumowaniu błędów dla $m_\mathcal{X}$ kolumn, otrzymujemy żądaną tezę.
\end{proof}

\subsubsection{Stabilność podprzestrzeni SVD i macierzowe nierówności koncentracyjne}

Ponieważ ekstrakcja właściwych kierunków z macierzy $\hat{X}$ opiera się na rozkładzie SVD, musimy zagwarantować stabilność wyznaczonych podprzestrzeni. Służy do tego klasyczne twierdzenie Wedina \cite{wedin1972perturbation} .

\begin{tw}[Twierdzenie Wedina o ograniczeniu perturbacji \cite{wedin1972perturbation}]\label{tw:wedin}
Jeśli istnieje $\bar{\alpha} > 0$ takie, że dla wartości osobliwych macierzy $B$ i $\hat{B}$ zachodzi separacja $ \min_{\hat{l}}|\sigma_{\hat{l}}(\hat{\Sigma}_1)| \ge \bar{\alpha} $ oraz znikają odpowiednie luki między blokami, to odległość podprzestrzeni rozpiętych przez pierwsze $k$ prawych wektorów osobliwych ($V_1$ i $\hat{V}_1$) można oszacować jako:
$$\|V_1 V_1^T - \hat{V}_1 \hat{V}_1^T\|_F \le \frac{2}{\bar{\alpha}}\|B - \hat{B}\|_F.$$

\end{tw}

Aplikując to do naszej macierzy i wykorzystując nierówność Weyla ($\sigma_k(\hat{X}^T) \ge \sigma_k(X^T) - \|X - \hat{X}\|_F$) , otrzymujemy:
$$\|V_1 V_1^T - \hat{V}_1 \hat{V}_1^T\|_F \le \frac{2\sqrt{m_\mathcal{X}}\nu_2}{\sigma_k(X^T) - \sqrt{m_\mathcal{X}}\nu_2}.$$
Aby to oszacowanie miało sens, $k$-ta wartość osobliwa oryginalnej macierzy $X^T$ musi być oddzielona od zera. Ponieważ $X^T X$ jest sumą losowych macierzy dodatnio półokreślonych, posłużymy się niezwykle silnym narzędziem, jakim jest macierzowa nierówność Chernoffa w sformułowaniu Troppa \cite{tropp2012user} .

\begin{tw}[Macierzowa nierówność Chernoffa \cite{tropp2012user}]\label{tw:matrixchernoff}
Rozważmy $X_1, \dots, X_m$ niezależnych, losowych macierzy dodatnio półokreślonych o wymiarze $k \times k$, spełniających prawie na pewno $\sigma_1(X_j) \le C$. Jeśli $\mu_{min} = \sigma_k(\sum_{j=1}^m \mathbb{E}X_j)$, to dla każdego $s \in (0, 1)$ zachodzi :
$$\mathbb{P}\left\{\sigma_k\left(\sum_{j=1}^m X_j\right) \le (1-s)\mu_{min}\right\} \le k e^{\frac{-\mu_{min} s^2}{2C}}.$$

\end{tw}

Prowadzi to nas do bezpośredniego oszacowania dla najmniejszej wartości osobliwej naszej macierzy:

\begin{lm}\label{lm:wartoscosobliwaX}
Dla dowolnego $s \in (0,1)$ zachodzi:
$$\sigma_k(X^T) \ge \sqrt{m_\mathcal{X}\alpha(1-s)}$$
z prawdopodobieństwem $1 - k e^{\frac{-m_\mathcal{X}\alpha s^2}{2kC_2^2}}$ .
\end{lm}

\begin{proof}
Dowód polega na zastosowaniu Twierdzenia \ref{tw:matrixchernoff} do macierzy kowariancji gradientów $\mathcal{G}^T \mathcal{G} = \sum_{j=1}^{m_\mathcal{X}} \nabla g(A\xi_j)\nabla g(A\xi_j)^T$ . Zauważamy, że $\sigma_1(\nabla g(A\xi_j)\nabla g(A\xi_j)^T) \le kC_2^2 =: C$ prawie na pewno . Wartość oczekiwana to $\mu_{min} = m_\mathcal{X}\sigma_k(H_g) \ge m_\mathcal{X}\alpha > 0$. Podstawiając te wartości do nierówności Chernoffa i biorąc pierwiastek kwadratowy, wprost otrzymujemy poszukiwane prawdopodobieństwo i ograniczenie .
\end{proof}

\subsubsection{Dowód Twierdzenia \ref{tw:zbieznoscalgorytmu2}}

Mając zabezpieczoną rekonstrukcję normy macierzy oraz stabilność jej podprzestrzeni, integrujemy te wyniki, by zakończyć dowód głównego twierdzenia dla wielowymiarowego algorytmu.

\begin{proof}
Korzystając łącznie z Wniosku \ref{wn:bladodzyskaniaX}, Twierdzenia \ref{tw:wedin} oraz Lematu \ref{lm:wartoscosobliwaX}, redukujemy czynnik $\sqrt{m_\mathcal{X}}$ z licznika i mianownika ograniczenia Wedina. Otrzymujemy, że z prawdopodobieństwem co najmniej
$$1 - \left(e^{-c_1' m_\Phi} + e^{-\sqrt{m_\Phi d}} + k e^{\frac{-m_\mathcal{X} \alpha s^2}{2kC_2^2}}\right)$$
zachodzi ograniczenie dla rzutów ortogonalnych (a z racji na konstrukcję algorytmu, również dla zrekonstruowanej macierzy $\hat{A}$) :
$$\|A^T A - \hat{A}^T \hat{A}\|_F = \|V_1 V_1^T - \hat{V}_1 \hat{V}_1^T\|_F \le \frac{2\nu_2}{\sqrt{\alpha(1-s)} - \nu_2}.$$

Przechodzimy teraz do wyznaczenia ostatecznego błędu na funkcji $f$. Ponieważ macierz $A$ posiada ortonormalne wiersze, zachodzi $A = A A^T A$. Wykorzystując to, dla dowolnego $x \in B_{\mathbb{R}^d}(1+\bar{\epsilon})$ mamy:
$$|f(x) - \hat{f}(x)| = |g(Ax) - g(A \hat{A}^T \hat{A} x)|$$
$$\le C_2\sqrt{k} \|Ax - A\hat{A}^T\hat{A}x\|_{l_2^k}$$
$$= C_2\sqrt{k} \|A(A^T A - \hat{A}^T \hat{A})x\|_{l_2^k}.$$
Korzystając z własności norm dla macierzy podrzędnej:
$$\le C_2\sqrt{k} \|(A^T A - \hat{A}^T \hat{A})\|_F \|x\|_{l_2^d}.$$
Ponieważ promień badanej kuli wynosi co najwyżej $1+\bar{\epsilon}$, wstawiamy uzyskane ograniczenie z Twierdzenia Wedina :
$$|f(x) - \hat{f}(x)| \le 2C_2\sqrt{k}(1+\bar{\epsilon})\frac{\nu_2}{\sqrt{\alpha(1-s)} - \nu_2},$$
co ostatecznie kończy dowód i potwierdza poprawność aproksymacji algorytmu wielowymiarowego. 
\end{proof}